\documentclass[10pt,a4paper]{amsart}
\usepackage{amsmath,amssymb,amsthm,graphicx}
% Retained preamble for the Stacks Project chapter snapshot in this directory.
% The chapter's own preamble.tex is a copy of the shared Stacks Project preamble; the block
% below is sliced MECHANICALLY, line for line, from it. Nothing was added, removed or
% redefined; the wrapper class, the xr-hyper/hyperref wiring and the externaldocument list
% are omitted because the assembler supplies the document class.
%
% lines taken from preamble.tex: 3, 5, 9, 11, 12, 156, 157, 158, ... (51 lines total)
%
% package coverage restored from the chapter's own preamble (2026-10-13): xy, xy-2cell

\usepackage[all]{xy}
\xyoption{2cell}
\UseAllTwocells

\newenvironment{reference}{\comment}{\endcomment}
\newenvironment{slogan}{\comment}{\endcomment}
\newenvironment{history}{\comment}{\endcomment}
\theoremstyle{plain}
\newtheorem{theorem}[subsection]{Theorem}
\newtheorem{proposition}[subsection]{Proposition}
\newtheorem{lemma}[subsection]{Lemma}
\theoremstyle{definition}
\newtheorem{definition}[subsection]{Definition}
\newtheorem{example}[subsection]{Example}
\newtheorem{exercise}[subsection]{Exercise}
\newtheorem{situation}[subsection]{Situation}
\theoremstyle{remark}
\newtheorem{remark}[subsection]{Remark}
\newtheorem{remarks}[subsection]{Remarks}
\numberwithin{equation}{subsection}
\def\lim{\mathop{\mathrm{lim}}\nolimits}
\def\colim{\mathop{\mathrm{colim}}\nolimits}
\def\Spec{\mathop{\mathrm{Spec}}}
\def\Hom{\mathop{\mathrm{Hom}}\nolimits}
\def\Ext{\mathop{\mathrm{Ext}}\nolimits}
\def\SheafHom{\mathop{\mathcal{H}\!\mathit{om}}\nolimits}
\def\SheafExt{\mathop{\mathcal{E}\!\mathit{xt}}\nolimits}
\def\Sch{\mathit{Sch}}
\def\Mor{\mathop{\mathrm{Mor}}\nolimits}
\def\Ob{\mathop{\mathrm{Ob}}\nolimits}
\def\Sh{\mathop{\mathit{Sh}}\nolimits}
\def\NL{\mathop{N\!L}\nolimits}
\def\CH{\mathop{\mathrm{CH}}\nolimits}
\def\proetale{{pro\text{-}\acute{e}tale}}
\def\etale{{\acute{e}tale}}
\def\QCoh{\mathit{QCoh}}
\def\Ker{\mathop{\mathrm{Ker}}}
\def\Im{\mathop{\mathrm{Im}}}
\def\Coker{\mathop{\mathrm{Coker}}}
\def\Coim{\mathop{\mathrm{Coim}}}
\DeclareMathSymbol{\boxtimes}{\mathbin}{AMSa}{"02}
\def\QCohstack{\mathcal{QC}\!\mathit{oh}}
\def\Cohstack{\mathcal{C}\!\mathit{oh}}
\def\Spacesstack{\mathcal{S}\!\mathit{paces}}
\def\Quotfunctor{\mathrm{Quot}}
\def\Hilbfunctor{\mathrm{Hilb}}
\def\Curvesstack{\mathcal{C}\!\mathit{urves}}
\def\Polarizedstack{\mathcal{P}\!\mathit{olarized}}
\def\Complexesstack{\mathcal{C}\!\mathit{omplexes}}
\def\Pic{\mathop{\mathrm{Pic}}\nolimits}
\def\Picardstack{\mathcal{P}\!\mathit{ic}}
\def\Picardfunctor{\mathrm{Pic}}
\def\Deformationcategory{\mathcal{D}\!\mathit{ef}}
% lane-s11: discard bibliographic \begin{reference} blocks entirely (incl. any \cite inside).
\newsavebox{\limrefbox}
\ifcsname reference\endcsname\expandafter\let\csname reference\endcsname\undefined\fi\ifcsname endreference\endcsname\expandafter\let\csname endreference\endcsname\undefined\fi
\def\citation#1{\relax}
\newenvironment{reference}{\setbox\limrefbox=\vbox\bgroup\let\cite\relax}{\egroup\ignorespaces}
\usepackage{cleveref}
\crefname{originalreference}{original reference}{original references}
\Crefname{originalreference}{Original reference}{Original references}
\makeatletter
\newlabel{morphisms-lemma-reach-points-scheme-theoretic-image}{{02JQ}{source}{}{sourcepointer-1}{}}
\newlabel{morphisms-lemma-reach-points-scheme-theoretic-image@cref}{{[originalreference][0][]02JQ}{[1][0][]source}{}{}{}}
\makeatother
\crefname{originalreference}{original reference}{original references}
\Crefname{originalreference}{Original reference}{Original references}
\makeatletter
\newlabel{schemes-section-points}{{01J5}{source}{}{sourcepointer-2}{}}
\newlabel{schemes-section-points@cref}{{[originalreference][0][]01J5}{[1][0][]source}{}{}{}}
\makeatother
\crefname{originalreference}{original reference}{original references}
\Crefname{originalreference}{Original reference}{Original references}
\makeatletter
\newlabel{algebra-lemma-exists-dvr}{{00PH}{source}{}{sourcepointer-3}{}}
\newlabel{algebra-lemma-exists-dvr@cref}{{[originalreference][0][]00PH}{[1][0][]source}{}{}{}}
\makeatother
\title{Reaching points and closures in the Noetherian limit case}
\author{The Stacks Project Authors}
\date{stacks-project:limits.tex:local-snapshot-20261012; GNU FDL 1.2 or later}
\begin{document}
\maketitle
\noindent\textit{Editorial note.} No omissions, hints or 'immediate' steps in the counted span.\par\medskip
\noindent\textit{Allowed standard prerequisites (editorial).}\begin{itemize}
\item Krull-Akizuki. Algebra, Krull-Akizuki lemma
\end{itemize}\medskip
\ifdefined\hypertarget\hypertarget{sourcepointer-1}{}\fi
\noindent\textit{Source reference.} Original reference 02JQ: sources/stacks-chow/tags.txt, source lines 3345--3345. Cross-chapter standard prerequisite cited by the author inside the counted span; it is used as a pointer and its own proof is not part of this item.\par\medskip
\ifdefined\hypertarget\hypertarget{sourcepointer-2}{}\fi
\noindent\textit{Source reference.} Original reference 01J5: sources/stacks-chow/tags.txt, source lines 2002--2002. Cross-chapter standard prerequisite cited by the author inside the counted span; it is used as a pointer and its own proof is not part of this item.\par\medskip
\ifdefined\hypertarget\hypertarget{sourcepointer-3}{}\fi
\noindent\textit{Source reference.} Original reference 00PH: sources/stacks-chow/tags.txt, source lines 931--931. Cross-chapter standard prerequisite cited by the author inside the counted span; it is used as a pointer and its own proof is not part of this item.\par\medskip
\begin{lemma}
\label{lemma-reach-point-closure-Noetherian}
Let $f : X \to Y$ be a morphism of schemes.
Assume $f$ finite type and $Y$ locally Noetherian.
Let $y \in Y$ be a point in the closure of the image of $f$.
Then there exists a commutative diagram
$$
\xymatrix{
\Spec(K) \ar[r] \ar[d] & X \ar[d]^f \\
\Spec(A) \ar[r] & Y
}
$$
where $A$ is a discrete valuation ring and $K$ is its field of fractions
mapping the closed point of $\Spec(A)$ to $y$. Moreover, we can assume
that the image point of $\Spec(K) \to X$ is a generic point $\eta$
of an irreducible component of $X$ and that $K = \kappa(\eta)$.
\end{lemma}

\begin{proof}
By the non-Noetherian version of this lemma
(Morphisms, Lemma \ref{morphisms-lemma-reach-points-scheme-theoretic-image})
there exists a point $x \in X$ such that $f(x)$ specializes to $y$.
We may replace $x$ by any point specializing to $x$, hence we may
assume that $x$ is a generic point of an irreducible component of $X$.
This produces a ring map $\mathcal{O}_{Y, y} \to \kappa(x)$
(see Schemes, Section \ref{schemes-section-points}).
Let $R \subset \kappa(x)$ be the image. Then $R$ is Noetherian as a quotient
of the Noetherian local ring $\mathcal{O}_{Y, y}$.
On the other hand, the extension $\kappa(x)$ is
a finitely generated extension of the fraction field of $R$
as $f$ is of finite type.
Thus there exists a discrete valuation ring $A \subset \kappa(x)$
with fraction field $\kappa(x)$ dominating $R$ by
Algebra, Lemma \ref{algebra-lemma-exists-dvr}. Then
$$
\xymatrix{
\Spec(\kappa(x)) \ar[d] \ar[rrr] & & & X \ar[d] \\
\Spec(A) \ar[r] & \Spec(R) \ar[r] & \Spec(\mathcal{O}_{Y, y}) \ar[r] & Y
}
$$
gives the desired diagram.
\end{proof}

\end{document}
